\newif\ifglobal

%%%%%%%%%%%%%%%%%%%%%%%
%%% Compile to view %%%
%%%%%%%%%%%%%%%%%%%%%%%

%%%%%%%%%%%%%%%%%%%%%%%%%%%%%%%%%%%%%%%%%%%%%%%%%%%
%%% Readme:                                     %%%
%%% https://www.overleaf.com/read/bwdjvfjksvjy/ %%%
%%%%%%%%%%%%%%%%%%%%%%%%%%%%%%%%%%%%%%%%%%%%%%%%%%%

% >>> M?: marks a module setting number or important notice

% >>> M4: Set {./relative-path} to external files for this module <<<
\def \modulefiles {./31-RMT}
% To add an external file, use:
% (Replace '---' with actual filename)
% '\input{\modulefiles/tables/---.tex}' for tables
% '\includegraphics{\modulefiles/figures/---}' for figures
% '\subfile{\modulefiles/---}' for register files

% >>> M1: This module is in global project? <<<
% 'YES' - uncomment; 'NO' - comment out
\globaltrue

\ifglobal
    \documentclass[main\_\_CN.tex]{subfiles}

\else
    % Font "FandolSong-Regular" used by ctex does not have GSUB table.
    % It causes warnings that do not affect compilation and can be ignored.
    \PassOptionsToPackage{quiet}{fontspec}

    \documentclass[openany, 10pt]{book}

    % >>> M2: In ./chip-spec-settings/preamble-trm-module, also find
    %         settings for visibility of todo notes, watermark <<<
    \usepackage{preamble-trm-module}


    % Enable Chinese typesetting
    \usepackage{ctex}

    % >>> M3: Temporary labels in Chinese
    %         you can add your own labels there <<<
    \myexternaldocument{temp-labels-cn}

\fi %\ifglobal

\setcounter{secnumdepth}{4}

% >>> M5: If this module needs any updates in
% './chip-spec-settings/preamble-trm-module.sty'
% (include additional package, etc.), contact your module PM
% or create the file './chip-spec-settings/changelog.tex' make the
% necessary updates yourself and document them in this file <<<

\raggedright


\begin{document}


% Sets language into which pre-defined bits of text are translated
\selectlanguage{chinese}

% Do not delete this block
\ifglobal
% Add the following front matter if
% this module is outside of global project
\else
    \InputIfFileExists{readme}{}{}
    {\let\clearpage\relax \listoftodos}
    {\let\clearpage\relax \listofdonetodos}
    \tableofcontents
    %\listoftables
    %\listoffigures
\fi


%%%%%%%%%%%%%%%%%%%%%%%%%
%%% TRM Module Begins %%%
%%%%%%%%%%%%%%%%%%%%%%%%%

\hypertarget{rmt}{}
\chapter{红外遥控 (RMT)}
\label{mod:rmt}

\section{概述}

RMT（红外遥控）是一个红外发送和接收控制器，可通过软件编解码多种红外协议。RMT 模块可以实现将模块内置 RAM 中的脉冲编码转换为信号输出，或检测模块的输入信号转换为脉冲编码存入 RAM 中。此外，RMT 模块可以选择是否对输出信号进行载波调制，也可以选择是否对输入信号进行滤波处理。


\section{功能描述}
\subsection{RMT 架构}

RMT 共有四个通道，可独立用于信号发送或接收，编码为 0 \verb+~+ 3，每个通道有一组功能相同的寄存器。为了方便叙述，以 \regindex{n} 表示各个通道。通道 0 默认从 block 0 读取发送数据或向 block 0 中存入接收数据；通道 1 默认从 block 1 读取发送数据或向 block 1 中存入接收数据，以此类推，如下图所示。

每个通道内部都有各自的一个时钟分频器，一个计数器，一个发射器和一个接收器。注意，同一个通道的发射器和接收器不可同时工作。四个通道共享一块 256 x 32 位的 RAM。

\begin{figure}[H]
\includegraphics[width=0.9\textwidth]{\modulefiles/figures/31-RMT-architecture-20211014.eps}
\centering
\caption{RMT 结构框图} \label{Figure:RMT2-1}
\end{figure}

\begin{figure}[H]
\includegraphics[width=0.9\textwidth]{\modulefiles/figures/31-RMT-channels-20211014.eps}
\centering
\caption{RMT 通道} \label{fig:rmt-channels}
\end{figure}

\subsection{RMT RAM}
\begin{figure}[htb!]
\includegraphics[width=1.0\textwidth]{\modulefiles/figures/rmt_data_structure.eps}
\centering
\caption{RAM中脉冲编码结构} \label{Figure:RMT2-2}
\end{figure}
RAM 中脉冲编码结构如图 \ref{Figure:RMT2-2} 所示。每个脉冲编码为 16 位，由 level 与 period 两部分组成。其中 level 表示输入或输出信号的逻辑电平值（0 或 1），period 表示该电平信号持续的时钟（图\ref{Figure:RMT2-1} clk\_div）周期数。Period 为 0 是传输结束标志。

RAM 按照 64 x 32 位分成四个 block。默认情况下每个通道只能使用一个 block（固定为通道 0 使用 block 0，通道 1 使用 block 1，以此类推）。当通道 \regindex{n} 单次发送或接收的脉冲编码数大于一个 block 时，可以进行乒乓操作，或者通过配置 \hyperref[fielddesc:RMTREGMEMSIZECHN]{RMT\_MEM\_SIZE\_CH\regindex{n}} 寄存器，允许该通道占用多个 block。当设置 \hyperref[fielddesc:RMTREGMEMSIZECHN]{RMT\_MEM\_SIZE\_CH\regindex{n}} > 1 时，通道 \regindex{n} 将占用 block (\regindex{n}) \verb+~+ block (\regindex{n} + \hyperref[fielddesc:RMTREGMEMSIZECHN]{RMT\_MEM\_SIZE\_CH\regindex{n}} -1) 的存储空间。通道 \regindex{n} +1 \verb+~+ \regindex{n} + \hyperref[fielddesc:RMTREGMEMSIZECHN]{RMT\_MEM\_SIZE\_CH\regindex{n}} - 1 因为对应的 RAM block 被占用而无法使用。注意，每个通道使用 RAM 的空间是根据地址从低到高进行映射的，因此通道 0 可以通过置位 \hyperref[fielddesc:RMTREGMEMSIZECHN]{RMT\_MEM\_SIZE\_CH\regindex{n}} 寄存器来使用通道 1、 2、 3 的 RAM 空间，但是通道 3 不能使用通道 0、 1 或 2 的 RAM 空间。

RAM 可被 APB 总线及通道的发射器或接收器访问，为了防止发射器和接收器访问 RAM 时发生冲突，用户可以通过配置 \hyperref[fielddesc:RMTREGMEMOWNERCHN]{RMT\_MEM\_OWNER\_CH\regindex{n}} 来决定当前 RAM 的使用权。当通道的发射器或接收器发生越权访问时会产生 \hyperref[fielddesc:RMTREGMEMOWNERERRCHN]{RMT\_MEM\_OWNER\_ERR\_CH\regindex{n}} 标志信号。

当 RMT 模块不工作时，可以通过配置 \hyperref[fielddesc:RMTREGMEMFORCEPD]{RMT\_MEM\_FORCE\_PD} 寄存器使 RAM 工作于低功耗模式。

\subsection{时钟}
用户可以通过配置 \hyperref[fielddesc:RMTREGREFALWAYSONCHN]{RMT\_REF\_ALWAYS\_ON\_CH\regindex{n}} 选择时钟分频器的驱动时钟：APB\_CLK 或者 REF\_TICK， 请参考\hyperref[mod:resclk]{复位和时钟}章节。\hyperref[fielddesc:RMTREGDIVCNTCHN]{RMT\_DIV\_CNT\_CH\regindex{n}} 寄存器可以配置分频器的分频系数，除 0 表示 256 分频外，其他分频数等同于寄存器 \hyperref[fielddesc:RMTREGDIVCNTCHN]{RMT\_DIV\_CNT\_CH\regindex{n}} 的值。时钟分频器可以通过配置 \hyperref[fielddesc:RMTREGREFCNTRSTCHN]{RMT\_REF\_CNT\_RST\_CH\regindex{n}} 进行复位。时钟分频器的分频时钟可供计数器使用。

\subsection{发射器}
当 \hyperref[fielddesc:RMTREGTXSTARTCHN]{RMT\_TX\_START\_CH\regindex{n}} 置为 1 时，通道 \regindex{n} 的发射器开始从通道对应 RAM block 的起始地址，按照地址从低到高依次读取脉冲编码进行发送。当遇到结束标志（period 等于 0）时，发射器将结束发送返回空闲状态，并产生 \hyperref[int:RMTCHnTXENDINT]{RMT\_CH\regindex{n}\_TX\_END\_INT} 中断。配置 \hyperref[fielddesc:RMTREGTXSTOPCHN]{RMT\_TX\_STOP\_CH\regindex{n}} 寄存器可以立刻停止发送并进入空闲状态。发射器空闲状态发送的电平由结束标志中的 level 段或者是 \hyperref[fielddesc:RMTREGIDLEOUTLVCHN]{RMT\_IDLE\_OUT\_LV\_CH\regindex{n}} 寄存器决定。用户可以配置 \hyperref[fielddesc:RMTREGIDLEOUTENCHN]{RMT\_IDLE\_OUT\_EN\_CH\regindex{n}} 寄存器来选择这两种方式。

当发送的脉冲编码较多时，可通过置位 \hyperref[fielddesc:RMTREGMEMTXWRAPEN]{RMT\_MEM\_TX\_WRAP\_EN} 寄存器使能乒乓操作。在乒乓操作模式下，发射器会循环从通道对应的 RAM 区域取出脉冲编码进行发送，直到遇到结束标识为止。例如，当 \hyperref[fielddesc:RMTREGMEMSIZECHN]{RMT\_MEM\_SIZE\_CH\regindex{n}} = 1 时，发射器将从 64 * \regindex{n} 地址开始发送，然后对应 RAM 的地址递增。发完 (64 * (\regindex{n} +1) - 1) 地址的数据后，下次继续从 64 * \regindex{n} 地址开始递增发送数据，依此类推，遇到结束标识时停止发送。\hyperref[fielddesc:RMTREGMEMSIZECHN]{RMT\_MEM\_SIZE\_CH\regindex{n}} > 1 的情形下，乒乓操作同样适用。

每当发射器发送的脉冲编码数大于等于 \hyperref[fielddesc:RMTREGTXLIMCHN]{RMT\_TX\_LIM\_CH\regindex{n}} 时，会产生 \hyperref[int:RMTCHnTXTHREVENTINT]{RMT\_CH\regindex{n}\_TX\_THR\_EVENT\_INT} 中断。在乒乓模式下，可以设置 \hyperref[fielddesc:RMTREGTXLIMCHN]{RMT\_TX\_LIM\_CH\regindex{n}} 为每个通道对应 RAM 空间的一半或几分之一。软件在检测到 \hyperref[int:RMTCHnTXTHREVENTINT]{RMT\_CH\regindex{n}\_TX\_THR\_EVENT\_INT} 中断之后，可以更新已使用过的 RAM 区域的脉冲编码，从而实现乒乓操作。

此外，发射器还可以对输出信号进行载波调制，置位 \hyperref[fielddesc:RMTREGCARRIERENCHN]{RMT\_CARRIER\_EN\_CH\regindex{n}} 可以使能该功能。载波的波形可配置。一个载波周期中高电平持续时间为 ( \hyperref[fielddesc:RMTREGCARRIERHIGHCHN]{RMT\_CARRIER\_HIGH\_CH\regindex{n}} +1) 个 APB\_CLK 或者 REF\_TICK 时钟周期，低电平持续的时间为 (\hyperref[fielddesc:RMTREGCARRIERLOWCHN]{RMT\_CARRIER\_LOW\_CH\regindex{n}} +1) 个 APB\_CLK 或者 REF\_TICK 时钟周期。置位 \hyperref[fielddesc:RMTREGCARRIEROUTLVCHN]{RMT\_CARRIER\_OUT\_LV\_CH\regindex{n}} 时在输出信号高电平上加载波信号，清零 \hyperref[fielddesc:RMTREGCARRIEROUTLVCHN]{RMT\_CARRIER\_OUT\_LV\_CH\regindex{n}} 时在输出信号低电平上加载波信号。同时，在进行载波调制时，载波可以一直加载在输出信号上，也可以仅加载在有效的脉冲编码（RAM 中的数据）上。通过配置 \hyperref[fielddesc:RMTREGCARRIEREFFENCHN]{RMT\_CARRIER\_EFF\_EN\_CH\regindex{n}} 寄存器，可以选择这两种模式。

置位 \hyperref[fielddesc:RMTREGTXCONTIMODECHN]{RMT\_TX\_CONTI\_MODE\_CH\regindex{n}} 寄存器可以使能发射器的持续发送功能。置位该寄存器后，发射器会循环发送 RAM 中的脉冲编码。配置 \hyperref[fielddesc:RMTREGTXLOOPCNTENCHN]{RMT\_TX\_LOOP\_CNT\_EN\_CH\regindex{n}} 后发射器会记录循环发送的次数。当该次数达到 \hyperref[fielddesc:RMTREGTXLOOPNUMCHN]{RMT\_TX\_LOOP\_NUM\_CH\regindex{n}} 设定的值时，会产生 \hyperref[int:RMTCHnTXLOOPINT]{RMT\_CH\regindex{n}\_TX\_LOOP\_INT} 中断。

置位 \hyperref[fielddesc:RMTREGTXSIMEN]{RMT\_TX\_SIM\_EN} 可以使能发射器多个通道同步发送的功能，此时多个通道会同时启动发送。\hyperref[fielddesc:RMTREGTXSIMCHN]{RMT\_TX\_SIM\_CH\regindex{n}} 用于选择同步发送的通道。

\subsection{接收器}
\hyperref[fielddesc:RMTREGRXENCHN]{RMT\_RX\_EN\_CH\regindex{n}} 置为 1 时接收器开始工作。接收器会检测信号电平及其持续的时钟周期数，将其按照脉冲编码的格式存入 RAM 中。当信号在一个电平下持续的时钟周期数超过 \hyperref[fielddesc:RMTREGIDLETHRESCHN]{RMT\_IDLE\_THRES\_CH\regindex{n}} 时，接收器结束接收过程，返回空闲状态，并产生 \hyperref[int:RMTCHnRXENDINT]{RMT\_CH\regindex{n}\_RX\_END\_INT} 中断。

每个通道都可以通过置位 \hyperref[fielddesc:RMTREGRXFILTERENCHN]{RMT\_RX\_FILTER\_EN\_CH\regindex{n}} 使能 接收器对输入信号进行滤波的功能。滤波器的功能为连续采样输入信号，如果输入信号在连续 RMT\_RX\_FILTER \_THRES\_CH\regindex{n} 个 APB  时钟周期内保持不变，则输入信号有效，否则输入信号无效。只有有效的输入信号才能通过滤波器。因此，滤波器会滤除脉冲宽度小于 \hyperref[fielddesc:RMTREGRXFILTERTHRESCHN]{RMT\_RX\_FILTER\_THRES\_CH\regindex{n}} 个 APB 时钟周期的线路毛刺。

\subsection{中断}
\begin{itemize}
\item \label{int:RMTCHnERRINT}RMT\_CH\regindex{n}\_ERR\_INT：当通道 \regindex{n}  发生读写数据不正确，或内存空满错误时，即触发此中断。
\item \label{int:RMTCHnTXTHREVENTINT}RMT\_CH\regindex{n}\_TX\_THR\_EVENT\_INT：发射器每发送 \hyperref[regdesc:RMTCHNTXLIMREG]{RMT\_CH\regindex{n}\_TX\_LIM\_REG} 的数据，即触发一次此中断。
\item \label{int:RMTCHnTXENDINT}RMT\_CH\regindex{n}\_TX\_END\_INT：当发射器停止发送信号时,即触发此中断。
\item \label{int:RMTCHnRXENDINT}RMT\_CH\regindex{n}\_RX\_END\_INT：当接收器停止接收信号时,即触发此中断。
\item \label{int:RMTCHnTXLOOPINT}RMT\_CH\regindex{n}\_TX\_LOOP\_INT：发射器处于循环发送模式时，当循环次数达到 \hyperref[fielddesc:RMTREGTXLOOPNUMCHN]{RMT\_TX\_LOOP\_NUM\_CH\regindex{n}} 的值后，会产生此中断。
\end{itemize}

\section{基地址} \label{RMT_BASE_ADDRESS}
用户可以通过两个不同的寄存器基地址访问 RMT，如表 \ref{table:rmt_base_address} 所示。更多信息，请访问章节 \ref{mod:sysmem} \textit{\nameref{mod:sysmem}}。

\begin{longtable}[c]{|l|C{2.5cm}|}
       \caption{RMT 基地址}
		\label{table:rmt_base_address} \\
	    \hline
        \rowcolor{lightgray}
        访问总线 & 基地址 \\
        \hline
        PeriBUS1 & \hexprefix{}3F416000 \\
        \hline
        PeriBUS2 & \hexprefix{}60016000 \\
        \hline
\end{longtable}

\hypertarget{rmt-reg-summ}{}
\section{寄存器列表}

请注意，下表中的地址都是相对于 RMT 基地址的地址偏移量（相对地址）。更多有关 RMT 基地址的信息，请前往第 \ref{RMT_BASE_ADDRESS} 节。

\begin{longtable}{ | p{6cm} | p{6cm} | p{2cm} | p{1cm} | }
\hline\rowcolor{lightgray}
名称 & 描述 & 地址 & 访问 \\ \hline
\endhead
\multicolumn{4}{|l|}{\textbf{配置寄存器}} \\ \hline
\hyperref[regdesc:RMTCHNCONF0REG]{RMT\_CH0CONF0\_REG} & 通道 0 配置寄存器 0 & \hexprefix{}0010 & 读/写 \\ \hline
\hyperref[regdesc:RMTCHNCONF1REG]{RMT\_CH0CONF1\_REG} & 通道 0 配置寄存器 1 & \hexprefix{}0014 & 不定 \\ \hline
\hyperref[regdesc:RMTCHNCONF0REG]{RMT\_CH1CONF0\_REG} & 通道 1 配置寄存器 0 & \hexprefix{}0018 & 读/写 \\ \hline
\hyperref[regdesc:RMTCHNCONF1REG]{RMT\_CH1CONF1\_REG} & 通道 1 配置寄存器 1 & \hexprefix{}001C & 不定 \\ \hline
\hyperref[regdesc:RMTCHNCONF0REG]{RMT\_CH2CONF0\_REG} & 通道 2 配置寄存器 0 & \hexprefix{}0020 & 读/写 \\ \hline
\hyperref[regdesc:RMTCHNCONF1REG]{RMT\_CH2CONF1\_REG} & 通道 2 配置寄存器 1 & \hexprefix{}0024 & 不定 \\ \hline
\hyperref[regdesc:RMTCHNCONF0REG]{RMT\_CH3CONF0\_REG} & 通道 3 配置寄存器 0 & \hexprefix{}0028 & 读/写 \\ \hline
\hyperref[regdesc:RMTCHNCONF1REG]{RMT\_CH3CONF1\_REG} & 通道 3 配置寄存器 1 & \hexprefix{}002C & 不定 \\ \hline
\hyperref[regdesc:RMTAPBCONFREG]{RMT\_APB\_CONF\_REG} & RMT APB 配置寄存器 & \hexprefix{}0080 & 读/写 \\ \hline
\hyperref[regdesc:RMTREFCNTRSTREG]{RMT\_REF\_CNT\_RST\_REG} & RMT 时钟分频器复位寄存器 & \hexprefix{}0088 & 读/写 \\ \hline
\hyperref[regdesc:RMTCHNRXCARRIERRMREG]{RMT\_CH0\_RX\_CARRIER\_RM\_REG} & 通道 0 载波清除寄存器 & \hexprefix{}008C & 读/写 \\ \hline
\hyperref[regdesc:RMTCHNRXCARRIERRMREG]{RMT\_CH1\_RX\_CARRIER\_RM\_REG} & 通道 1 载波清除寄存器 & \hexprefix{}0090 & 读/写 \\ \hline
\hyperref[regdesc:RMTCHNRXCARRIERRMREG]{RMT\_CH2\_RX\_CARRIER\_RM\_REG} & 通道 2 载波清除寄存器 & \hexprefix{}0094 & 读/写 \\ \hline
\hyperref[regdesc:RMTCHNRXCARRIERRMREG]{RMT\_CH3\_RX\_CARRIER\_RM\_REG} & 通道 3 载波清除寄存器 & \hexprefix{}0098 & 读/写 \\ \hline
\multicolumn{4}{|l|}{\textbf{载波占空比寄存器}} \\ \hline
\hyperref[regdesc:RMTCHNCARRIERDUTYREG]{RMT\_CH0CARRIER\_DUTY\_REG} & 通道 0 占空比配置寄存器 & \hexprefix{}0060 & 读/写 \\ \hline
\hyperref[regdesc:RMTCHNCARRIERDUTYREG]{RMT\_CH1CARRIER\_DUTY\_REG} & 通道 1 占空比配置寄存器 & \hexprefix{}0064 & 读/写 \\ \hline
\hyperref[regdesc:RMTCHNCARRIERDUTYREG]{RMT\_CH2CARRIER\_DUTY\_REG} & 通道 2 占空比配置寄存器 & \hexprefix{}0068 & 读/写 \\ \hline
\hyperref[regdesc:RMTCHNCARRIERDUTYREG]{RMT\_CH3CARRIER\_DUTY\_REG} & 通道 3 占空比配置寄存器 & \hexprefix{}006C & 读/写 \\ \hline
\multicolumn{4}{|l|}{\textbf{发送事件配置寄存器}} \\ \hline
\hyperref[regdesc:RMTCHNTXLIMREG]{RMT\_CH0\_TX\_LIM\_REG} & 通道 0 Tx 配置寄存器 & \hexprefix{}0070 & 不定 \\ \hline
\hyperref[regdesc:RMTCHNTXLIMREG]{RMT\_CH1\_TX\_LIM\_REG} & 通道 1 Tx 配置寄存器 & \hexprefix{}0074 & 不定 \\ \hline
\hyperref[regdesc:RMTCHNTXLIMREG]{RMT\_CH2\_TX\_LIM\_REG} & 通道 2 Tx 配置寄存器 & \hexprefix{}0078 & 不定 \\ \hline
\hyperref[regdesc:RMTCHNTXLIMREG]{RMT\_CH3\_TX\_LIM\_REG} & 通道 3 Tx 配置寄存器 & \hexprefix{}007C & 不定 \\ \hline
\hyperref[regdesc:RMTTXSIMREG]{RMT\_TX\_SIM\_REG} & RMT 同步发送寄存器 & \hexprefix{}0084 & 读/写 \\ \hline
\multicolumn{4}{|l|}{\textbf{状态寄存器}} \\ \hline
\hyperref[regdesc:RMTCHNSTATUSREG]{RMT\_CH0STATUS\_REG} & 通道 0 状态寄存器 & \hexprefix{}0030 & 只读 \\ \hline
\hyperref[regdesc:RMTCHNSTATUSREG]{RMT\_CH1STATUS\_REG} & 通道 1 状态寄存器 & \hexprefix{}0034 & 只读 \\ \hline
\hyperref[regdesc:RMTCHNSTATUSREG]{RMT\_CH2STATUS\_REG} & 通道 2 状态寄存器 & \hexprefix{}0038 & 只读 \\ \hline
\hyperref[regdesc:RMTCHNSTATUSREG]{RMT\_CH3STATUS\_REG} & 通道 3 状态寄存器 & \hexprefix{}003C & 只读 \\ \hline
\hyperref[regdesc:RMTCHNADDRREG]{RMT\_CH0ADDR\_REG} & 通道 0 地址寄存器 & \hexprefix{}0040 & 只读 \\ \hline
\hyperref[regdesc:RMTCHNADDRREG]{RMT\_CH1ADDR\_REG} & 通道 1 地址寄存器 & \hexprefix{}0044 & 只读 \\ \hline
\hyperref[regdesc:RMTCHNADDRREG]{RMT\_CH2ADDR\_REG} & 通道 2 地址寄存器 & \hexprefix{}0048 & 只读 \\ \hline
\hyperref[regdesc:RMTCHNADDRREG]{RMT\_CH3ADDR\_REG} & 通道 3 地址寄存器 & \hexprefix{}004C & 只读 \\ \hline
\multicolumn{4}{|l|}{\textbf{版本寄存器}} \\ \hline
\hyperref[regdesc:RMTDATEREG]{RMT\_DATE\_REG} & 版本控制寄存器 & \hexprefix{}00FC & 读/写 \\ \hline
\multicolumn{4}{|l|}{\textbf{FIFO 读/写 寄存器}} \\ \hline
\hyperref[regdesc:RMTCHNDATAREG]{RMT\_CH0DATA\_REG} & 配置 APB FIFO 对通道 0 进行读写操作 & \hexprefix{}0000 & 只读 \\ \hline
\hyperref[regdesc:RMTCHNDATAREG]{RMT\_CH1DATA\_REG} & 配置 APB FIFO 对通道 1 进行读写操作 & \hexprefix{}0004 & 只读 \\ \hline
\hyperref[regdesc:RMTCHNDATAREG]{RMT\_CH2DATA\_REG} & 配置 APB FIFO 对通道 2 进行读写操作 & \hexprefix{}0008 & 只读 \\ \hline
\hyperref[regdesc:RMTCHNDATAREG]{RMT\_CH3DATA\_REG} & 配置 APB FIFO 对通道 3 进行读写操作 & \hexprefix{}000C & 只读 \\ \hline
\multicolumn{4}{|l|}{\textbf{中断寄存器}} \\ \hline
\hyperref[regdesc:RMTINTRAWREG]{RMT\_INT\_RAW\_REG} & 原始中断状态寄存器 & \hexprefix{}0050 & 只读 \\ \hline
\hyperref[regdesc:RMTINTSTREG]{RMT\_INT\_ST\_REG} & 隐蔽中断状态寄存器 & \hexprefix{}0054 & 只读 \\ \hline
\hyperref[regdesc:RMTINTENAREG]{RMT\_INT\_ENA\_REG} & 中断使能寄存器 & \hexprefix{}0058 & 读/写 \\ \hline
\hyperref[regdesc:RMTINTCLRREG]{RMT\_INT\_CLR\_REG} & 中断清除寄存器 & \hexprefix{}005C & 只写 \\ \hline
\end{longtable}


\section{寄存器}

\begin{register}{H}{RMT\_CH\regindex{n}CONF0\_REG (\regindex{n}: 0-3)}{\hexprefix{}0010+8*\regindex{n}}\label{regdesc:RMTCHNCONF0REG}
\regfield{(reserved)}{2}{30}{00}%
\regfield{RMT\_CARRIER\_OUT\_LV\_CH\regindex{n}}{1}{29}{{1}}%
\regfield{RMT\_CARRIER\_EN\_CH\regindex{n}}{1}{28}{{1}}%
\regfield{RMT\_CARRIER\_EFF\_EN\_CH\regindex{n}}{1}{27}{{1}}%
\regfield{RMT\_MEM\_SIZE\_CH\regindex{n}}{3}{24}{{\hexprefix{}1}}%
\regfield{RMT\_IDLE\_THRES\_CH\regindex{n}}{16}{8}{{\hexprefix{}1000}}%
\regfield{RMT\_DIV\_CNT\_CH\regindex{n}}{8}{0}{{\hexprefix{}2}}%
\reglabel{Reset}\regnewline%
\begin{regdesc}\begin{reglist}
\label{fielddesc:RMTREGDIVCNTCHN}\item [RMT\_DIV\_CNT\_CH\regindex{n}] 用于配置通道 \regindex{n} 分频器的分频系数。（读/写）
\label{fielddesc:RMTREGIDLETHRESCHN}\item [RMT\_IDLE\_THRES\_CH\regindex{n}] 当接收器长时间检测不到信号沿变化，即计数器的值大于等于此寄存器的值时，接收器结束接收过程。（读/写）
\label{fielddesc:RMTREGMEMSIZECHN}\item [RMT\_MEM\_SIZE\_CH\regindex{n}] 配置通道 \regindex{n} 可以使用的 block 数量，范围为 1 \verb+~+ 4-\regindex{n}。（读/写）
\label{fielddesc:RMTREGCARRIEREFFENCHN}\item [RMT\_CARRIER\_EFF\_EN\_CH\regindex{n}] 1：配置通道 \regindex{n} 仅在发送状态下对输出信号载波调制；0：配置通道 \regindex{n} 对所有状态（空闲状态、从 RAM 中读取数据和发送数据阶段）均加载载波。仅在 \hyperref[fielddesc:RMTREGCARRIERENCHN]{RMT\_CARRIER\_EN\_CH\regindex{n}} 为 1 时有效。（读/写）
\label{fielddesc:RMTREGCARRIERENCHN}\item [RMT\_CARRIER\_EN\_CH\regindex{n}] 通道 \regindex{n} 的载波调制使能位。 1：对输出信号进行载波调制；0：禁止对输出信号进行载波调制。（读/写）
\label{fielddesc:RMTREGCARRIEROUTLVCHN}\item [RMT\_CARRIER\_OUT\_LV\_CH\regindex{n}] 用于配置通道 \regindex{n} 的载波调制方式。
1'h0：载波加载在输出信号低电平上；
1'h1：载波加载在输出信号高电平上。（读/写）
\end{reglist}\end{regdesc}
\end{register}


\begin{register}{H}{RMT\_CH\regindex{n}CONF1\_REG (\regindex{n}: 0-3)}{\hexprefix{}0014+8*\regindex{n}}\label{regdesc:RMTCHNCONF1REG}
\regfield{(reserved)}{11}{21}{00000000000}%
\regfield{RMT\_TX\_STOP\_CH\regindex{n}}{1}{20}{{0}}%
\regfield{RMT\_IDLE\_OUT\_EN\_CH\regindex{n}}{1}{19}{{0}}%
\regfield{RMT\_IDLE\_OUT\_LV\_CH\regindex{n}}{1}{18}{{0}}%
\regfield{RMT\_REF\_ALWAYS\_ON\_CH\regindex{n}}{1}{17}{{0}}%
\regfield{RMT\_CHK\_RX\_CARRIER\_EN\_CH\regindex{n}}{1}{16}{{0}}%
\regfield{RMT\_RX\_FILTER\_THRES\_CH\regindex{n}}{8}{8}{{\hexprefix{}f}}%
\regfield{RMT\_RX\_FILTER\_EN\_CH\regindex{n}}{1}{7}{{0}}%
\regfield{RMT\_TX\_CONTI\_MODE\_CH\regindex{n}}{1}{6}{{0}}%
\regfield{RMT\_MEM\_OWNER\_CH\regindex{n}}{1}{5}{{1}}%
\regfield{RMT\_APB\_MEM\_RST\_CH\regindex{n}}{1}{4}{{0}}%
\regfield{RMT\_MEM\_RD\_RST\_CH\regindex{n}}{1}{3}{{0}}%
\regfield{RMT\_MEM\_WR\_RST\_CH\regindex{n}}{1}{2}{{0}}%
\regfield{RMT\_RX\_EN\_CH\regindex{n}}{1}{1}{{0}}%
\regfield{RMT\_TX\_START\_CH\regindex{n}}{1}{0}{{0}}%
\reglabel{Reset}\regnewline%
\begin{regdesc}\begin{reglist}
\label{fielddesc:RMTREGTXSTARTCHN}\item [RMT\_TX\_START\_CH\regindex{n}] 用于使能通道 \regindex{n} 的发射器。发射器开始发送数据。（读/写）
\label{fielddesc:RMTREGRXENCHN}\item [RMT\_RX\_EN\_CH\regindex{n}] 用于使能通道 \regindex{n} 的接收器。接收器开始接收数据。（读/写）
\label{fielddesc:RMTREGMEMWRRSTCHN}\item [RMT\_MEM\_WR\_RST\_CH\regindex{n}] 用于复位通道 \regindex{n} 中接收器访问的 RAM 写地址。（只写）
\label{fielddesc:RMTREGMEMRDRSTCHN}\item [RMT\_MEM\_RD\_RST\_CH\regindex{n}] 用于复位通道 \regindex{n} 中发射器访问的 RAM 读地址。（只写）
\label{fielddesc:RMTREGMEMOWNERCHN}\item [RMT\_MEM\_OWNER\_CH\regindex{n}] 标志通道 \regindex{n} 的 RAM 使用权。
1'h1：接收器有权使用该通道的 RAM；
1'h0：发射器有权使用该通道的 RAM。（读/写）
\label{fielddesc:RMTREGTXCONTIMODECHN}\item [RMT\_TX\_CONTI\_MODE\_CH\regindex{n}] 设置通道 \regindex{n} 发送结束后，发射器不进入空闲状态，而是继续从第一个字节开始发送数据。（读/写）
\label{fielddesc:RMTREGRXFILTERENCHN}\item [RMT\_RX\_FILTER\_EN\_CH\regindex{n}] 使能通道 \regindex{n} 接收器的滤波功能。（读/写）
\label{fielddesc:RMTREGRXFILTERTHRESCHN}\item [RMT\_RX\_FILTER\_THRES\_CH\regindex{n}] 接收模式下，忽略宽度小于 RMT\_RX\_FILTER\_THRES\_CH\regindex{n} 个 APB 时钟周期的脉冲。（读/写）
\label{fielddesc:RMTREGCHKRXCARRIERENCHN}\item [RMT\_CHK\_RX\_CARRIER\_EN\_CH\regindex{n}] 用于使能 RAM 读循环模式（通道 \regindex{n} 中需使能载波调制模式）。（读/写）
\label{fielddesc:RMTREGREFALWAYSONCHN}\item [RMT\_REF\_ALWAYS\_ON\_CH\regindex{n}] 用于选择通道 \regindex{n} 的基础时钟。
1'h1：APB\_CLK；1'h0：REF\_TICK（读/写）
\label{fielddesc:RMTREGIDLEOUTLVCHN}\item [RMT\_IDLE\_OUT\_LV\_CH\regindex{n}] 配置通道 \regindex{n} 处于空闲状态下的输出信号电平。（读/写）
\label{fielddesc:RMTREGIDLEOUTENCHN}\item [RMT\_IDLE\_OUT\_EN\_CH\regindex{n}] 通道 \regindex{n} 处于空闲状态下的输出使能位。（读/写）
\label{fielddesc:RMTREGTXSTOPCHN}\item [RMT\_TX\_STOP\_CH\regindex{n}] 通道 \regindex{n} 中发射器停止发送信息。（读/写）
\end{reglist}\end{regdesc}
\end{register}


\begin{register}{H}{RMT\_APB\_CONF\_REG}{\hexprefix{}0080}\label{regdesc:RMTAPBCONFREG}
\regfield{RMT\_CLK\_EN}{1}{31}{{0}}%
\regfield{(reserved)}{26}{5}{00000000000000000000000000}%
\regfield{RMT\_MEM\_FORCE\_PU}{1}{4}{{0}}%
\regfield{RMT\_MEM\_FORCE\_PD}{1}{3}{{0}}%
\regfield{RMT\_MEM\_CLK\_FORCE\_ON}{1}{2}{{1}}%
\regfield{RMT\_MEM\_TX\_WRAP\_EN}{1}{1}{{0}}%
\regfield{RMT\_APB\_FIFO\_MASK}{1}{0}{{0}}%
\reglabel{Reset}\regnewline%
\begin{regdesc}\begin{reglist}
\label{fielddesc:RMTREGAPBFIFOMASK}\item [RMT\_APB\_FIFO\_MASK] 1'h1：直接访问 RAM；1'h0：经 APB FIFO 访问 RAM。（读/写）
\label{fielddesc:RMTREGMEMTXWRAPEN}\item [RMT\_MEM\_TX\_WRAP\_EN] 乒乓模式使能位。该模式下，当发射器发送的数据量大于 block 大小，发射器将继续从第一字节开始发送数据。（读/写）
\label{fielddesc:RMTREGMEMCLKFORCEON}\item [RMT\_MEM\_CLK\_FORCE\_ON] RMT 模块工作时，使能 RAM 的时钟；RMT 模块不工作时，关闭 RAM 的时钟，从而实现低功耗。（读/写）
\label{fielddesc:RMTREGMEMFORCEPD}\item [RMT\_MEM\_FORCE\_PD] 降低 RMT RAM 的功耗。（读/写）
\label{fielddesc:RMTREGMEMFORCEPU}\item [RMT\_MEM\_FORCE\_PU] 1：禁用 RMT RAM 的 Light-sleep 低功耗功能；0：RMT 处于 Light-sleep 模式时，设置 RAM 进入低功耗模式。（读/写）
\label{fielddesc:RMTREGCLKEN}\item [RMT\_CLK\_EN] RMT 寄存器的时钟门控使能位，用于 RMT 寄存器低功耗控制。1：使能 RMT 寄存器驱动时钟；0：禁用 RMT 寄存器驱动时钟。（读/写）
\end{reglist}\end{regdesc}
\end{register}


\begin{register}{H}{RMT\_REF\_CNT\_RST\_REG}{\hexprefix{}0088}\label{regdesc:RMTREFCNTRSTREG}
\regfield{(reserved)}{28}{4}{0000000000000000000000000000}%
\regfield{RMT\_REF\_CNT\_RST\_CH3}{1}{3}{{0}}%
\regfield{RMT\_REF\_CNT\_RST\_CH2}{1}{2}{{0}}%
\regfield{RMT\_REF\_CNT\_RST\_CH1}{1}{1}{{0}}%
\regfield{RMT\_REF\_CNT\_RST\_CH0}{1}{0}{{0}}%
\reglabel{Reset}\regnewline%
\begin{regdesc}\begin{reglist}
\label{fielddesc:RMTREGREFCNTRSTCHN}\item [RMT\_REF\_CNT\_RST\_CH\regindex{n}] 复位通道 \regindex{n} 的时钟分频器。（读/写）
\end{reglist}\end{regdesc}
\end{register}


\begin{register}{H}{RMT\_CH\regindex{n}\_RX\_CARRIER\_RM\_REG (\regindex{n}: 0-3)}{\hexprefix{}008C+4*\regindex{n}}\label{regdesc:RMTCHNRXCARRIERRMREG}
\regfield{RMT\_CARRIER\_HIGH\_THRES\_CH\regindex{n}}{16}{16}{{\hexprefix{}00}}%
\regfield{RMT\_CARRIER\_LOW\_THRES\_CH\regindex{n}}{16}{0}{{\hexprefix{}00}}%
\reglabel{Reset}\regnewline%
\begin{regdesc}\begin{reglist}
\label{fielddesc:RMTREGCARRIERLOWTHRESCHN}\item [RMT\_CARRIER\_LOW\_THRES\_CH\regindex{n}] 载波调制模式下，通道 \regindex{n} 低电平周期为 (RMT\_CARRIER\_LOW\_THRES\_CH\regindex{n} + 1) 个时钟周期。（读/写）
\label{fielddesc:RMTREGCARRIERHIGHTHRESCHN}\item [RMT\_CARRIER\_HIGH\_THRES\_CH\regindex{n}] 载波模式下，通道 \regindex{n} 高电平周期为 (RMT\_CARRIER\_HIGH\_THRES\_CH\regindex{n} + 1) 个时钟周期。（读/写）
\end{reglist}\end{regdesc}
\end{register}


\begin{register}{H}{RMT\_CH\regindex{n}CARRIER\_DUTY\_REG (\regindex{n}: 0-3)}{\hexprefix{}0060+4*\regindex{n}}\label{regdesc:RMTCHNCARRIERDUTYREG}
\regfield{RMT\_CARRIER\_HIGH\_CH\regindex{n}}{16}{16}{{\hexprefix{}40}}%
\regfield{RMT\_CARRIER\_LOW\_CH\regindex{n}}{16}{0}{{\hexprefix{}40}}%
\reglabel{Reset}\regnewline%
\begin{regdesc}\begin{reglist}
\label{fielddesc:RMTREGCARRIERLOWCHN}\item [RMT\_CARRIER\_LOW\_CH\regindex{n}] 用于配置通道 \regindex{n} 载波的低电平时钟周期。（读/写）
\label{fielddesc:RMTREGCARRIERHIGHCHN}\item [RMT\_CARRIER\_HIGH\_CH\regindex{n}] 用于配置通道 \regindex{n} 载波的高电平时钟周期。（读/写）
\end{reglist}\end{regdesc}
\end{register}


\begin{register}{H}{RMT\_CH\regindex{n}\_TX\_LIM\_REG (\regindex{n}: 0-3)}{\hexprefix{}0070+4*\regindex{n}}\label{regdesc:RMTCHNTXLIMREG}
\regfield{(reserved)}{11}{21}{00000000000}%
\regfield{RMT\_LOOP\_COUNT\_RESET\_CH\regindex{n}}{1}{20}{{0}}%
\regfield{RMT\_TX\_LOOP\_CNT\_EN\_CH\regindex{n}}{1}{19}{{0}}%
\regfield{RMT\_TX\_LOOP\_NUM\_CH\regindex{n}}{10}{9}{{\hexprefix{}0}}%
\regfield{RMT\_TX\_LIM\_CH\regindex{n}}{9}{0}{{\hexprefix{}80}}%
\reglabel{Reset}\regnewline%
\begin{regdesc}\begin{reglist}
\label{fielddesc:RMTREGTXLIMCHN}\item [RMT\_TX\_LIM\_CH\regindex{n}] 用于配置通道 \regindex{n} 发送脉冲编码数量的上限值。当 \hyperref[fielddesc:RMTREGMEMSIZECHN]{RMT\_MEM\_SIZE\_CH\regindex{n}} = 1 时，此寄存器可设置为 0 \verb+~+ 128 (64 * 32 /16 = 128) 之间的值；当 \hyperref[fielddesc:RMTREGMEMSIZECHN]{RMT\_MEM\_SIZE\_CH\regindex{n}} > 1 时，此寄存器可设置为 (0 \verb+~+ 128) * \hyperref[fielddesc:RMTREGMEMSIZECHN]{RMT\_MEM\_SIZE\_CH\regindex{n}} 之间的值。（读/写）
\label{fielddesc:RMTREGTXLOOPNUMCHN}\item [RMT\_TX\_LOOP\_NUM\_CH\regindex{n}] 用于配置连续发送模式下最大循环发送次数。（读/写）
\label{fielddesc:RMTREGTXLOOPCNTENCHN}\item [RMT\_TX\_LOOP\_CNT\_EN\_CH\regindex{n}] 使能循环计数。（读/写）
\label{fielddesc:RMTREGLOOPCOUNTRESETCHN}\item [RMT\_LOOP\_COUNT\_RESET\_CH\regindex{n}] 重置 RMT\_TX\_CONTI\_MODE\_CH\regindex{n} 模式下循环计数器。（只写）
\end{reglist}\end{regdesc}
\end{register}


\begin{register}{H}{RMT\_TX\_SIM\_REG}{\hexprefix{}0084}\label{regdesc:RMTTXSIMREG}
\regfield{(reserved)}{27}{5}{000000000000000000000000000}%
\regfield{RMT\_TX\_SIM\_EN}{1}{4}{{0}}%
\regfield{RMT\_TX\_SIM\_CH3}{1}{3}{{0}}%
\regfield{RMT\_TX\_SIM\_CH2}{1}{2}{{0}}%
\regfield{RMT\_TX\_SIM\_CH1}{1}{1}{{0}}%
\regfield{RMT\_TX\_SIM\_CH0}{1}{0}{{0}}%
\reglabel{Reset}\regnewline%
\begin{regdesc}\begin{reglist}
\label{fielddesc:RMTREGTXSIMCHN}\item [RMT\_TX\_SIM\_CH\regindex{n}] 使能通道 \regindex{n} 与其他启用的通道同步开始发送数据。（读/写）
\label{fielddesc:RMTREGTXSIMEN}\item [RMT\_TX\_SIM\_EN] 使能多通道同步发送数据。（读/写）
\end{reglist}\end{regdesc}
\end{register}


\begin{register}{H}{RMT\_CH\regindex{n}STATUS\_REG (\regindex{n}: 0-3)}{\hexprefix{}0030+4*\regindex{n}}\label{regdesc:RMTCHNSTATUSREG}
\regfield{(reserved)}{4}{28}{0000}%
\regfield{RMT\_APB\_MEM\_RD\_ERR\_CH\regindex{n}}{1}{27}{{0}}%
\regfield{RMT\_APB\_MEM\_WR\_ERR\_CH\regindex{n}}{1}{26}{{0}}%
\regfield{RMT\_MEM\_EMPTY\_CH\regindex{n}}{1}{25}{{0}}%
\regfield{RMT\_MEM\_FULL\_CH\regindex{n}}{1}{24}{{0}}%
\regfield{RMT\_MEM\_OWNER\_ERR\_CH\regindex{n}}{1}{23}{{0}}%
\regfield{RMT\_STATE\_CH\regindex{n}}{3}{20}{{\hexprefix{}0}}%
\regfield{(reserved)}{1}{19}{0}%
\regfield{RMT\_MEM\_RADDR\_EX\_CH\regindex{n}}{9}{10}{{\hexprefix{}0}}%
\regfield{(reserved)}{1}{9}{0}%
\regfield{RMT\_MEM\_WADDR\_EX\_CH\regindex{n}}{9}{0}{{\hexprefix{}0}}%
\reglabel{Reset}\regnewline%
\begin{regdesc}\begin{reglist}
\label{fielddesc:RMTREGMEMWADDREXCHN}\item [RMT\_MEM\_WADDR\_EX\_CH\regindex{n}] 记录通道 \regindex{n} 中接收器使用 RAM 时的地址偏移量。（只读）
\label{fielddesc:RMTREGMEMRADDREXCHN}\item [RMT\_MEM\_RADDR\_EX\_CH\regindex{n}] 记录通道 \regindex{n} 中发射器使用 RAM 时的地址偏移量。（只读）
\label{fielddesc:RMTREGSTATECHN}\item [RMT\_STATE\_CH\regindex{n}] 记录通道 \regindex{n} 的 FSM 状态。（只读）
\label{fielddesc:RMTREGMEMOWNERERRCHN}\item [RMT\_MEM\_OWNER\_ERR\_CH\regindex{n}] RAM block 使用权发生错误时，将触发此状态位。（只读）
\label{fielddesc:RMTREGMEMFULLCHN}\item [RMT\_MEM\_FULL\_CH\regindex{n}] 接收器接收的数据量大于 RAM block 时，将触发此状态位。（只读）
\label{fielddesc:RMTREGMEMEMPTYCHN}\item [RMT\_MEM\_EMPTY\_CH\regindex{n}] 待发送的数据量大于 RAM block，且乒乓模式未启用时，将触发此状态位。（只读）
\label{fielddesc:RMTREGAPBMEMWRERRCHN}\item [RMT\_APB\_MEM\_WR\_ERR\_CH\regindex{n}] 通过 APB 总线执行写操作时，如果偏移地址溢出 RAM block，将触发此状态位。（只读）
\label{fielddesc:RMTREGAPBMEMRDERRCHN}\item [RMT\_APB\_MEM\_RD\_ERR\_CH\regindex{n}] 通过 APB 总线执行读操作时，如果偏移地址溢出 RAM block，将触发此状态位。（只读）
\end{reglist}\end{regdesc}
\end{register}


\begin{register}{H}{RMT\_CH\regindex{n}ADDR\_REG (\regindex{n}: 0-3)}{\hexprefix{}0040+4*\regindex{n}}\label{regdesc:RMTCHNADDRREG}
\regfield{(reserved)}{13}{19}{0000000000000}%
\regfield{RMT\_APB\_MEM\_RADDR\_CH\regindex{n}}{9}{10}{{\hexprefix{}0}}%
\regfield{(reserved)}{1}{9}{0}%
\regfield{RMT\_APB\_MEM\_WADDR\_CH\regindex{n}}{9}{0}{{\hexprefix{}0}}%
\reglabel{Reset}\regnewline%
\begin{regdesc}\begin{reglist}
\label{fielddesc:RMTREGAPBMEMWADDRCHN}\item [RMT\_APB\_MEM\_WADDR\_CH\regindex{n}] 记录通过 APB 总线对 RAM 执行写操作时，地址的偏移量。（只读）
\label{fielddesc:RMTREGAPBMEMRADDRCHN}\item [RMT\_APB\_MEM\_RADDR\_CH\regindex{n}] 记录通过 APB 总线对 RAM 执行读操作时，地址的偏移量。（只读）
\end{reglist}\end{regdesc}
\end{register}


\begin{register}{H}{RMT\_DATE\_REG}{\hexprefix{}00FC}\label{regdesc:RMTDATEREG}
\regfield{RMT\_DATE}{32}{0}{{\hexprefix{}19072601}}%
\reglabel{Reset}\regnewline%
\begin{regdesc}\begin{reglist}
\label{fielddesc:RMTREGDATE}\item [RMT\_DATE] 版本控制寄存器（读/写）
\end{reglist}\end{regdesc}
\end{register}


\begin{register}{H}{RMT\_CH\regindex{n}DATA\_REG (\regindex{n}: 0-3)}{\hexprefix{}0000+4*\regindex{n}}\label{regdesc:RMTCHNDATAREG}
\regfield{RMT\_CH\regindex{n}\_DATA\_REG}{32}{0}{{\hexprefix{}000000}}%
\reglabel{Reset}\regnewline%
\begin{regdesc}\begin{reglist}
\label{fielddesc:RMTCHNDATAREG}\item [RMT\_CH\regindex{n}\_DATA\_REG] 通过 APB FIFO 对通道 \regindex{n} 进行读写操作。（只读）
\end{reglist}\end{regdesc}
\end{register}


\begin{register}{H}{RMT\_INT\_RAW\_REG}{\hexprefix{}0050}\label{regdesc:RMTINTRAWREG}
\regfield{(reserved)}{12}{20}{000000000000}%
\regfield{RMT\_CH3\_TX\_LOOP\_INT\_RAW}{1}{19}{{0}}%
\regfield{RMT\_CH2\_TX\_LOOP\_INT\_RAW}{1}{18}{{0}}%
\regfield{RMT\_CH1\_TX\_LOOP\_INT\_RAW}{1}{17}{{0}}%
\regfield{RMT\_CH0\_TX\_LOOP\_INT\_RAW}{1}{16}{{0}}%
\regfield{RMT\_CH3\_TX\_THR\_EVENT\_INT\_RAW}{1}{15}{{0}}%
\regfield{RMT\_CH2\_TX\_THR\_EVENT\_INT\_RAW}{1}{14}{{0}}%
\regfield{RMT\_CH1\_TX\_THR\_EVENT\_INT\_RAW}{1}{13}{{0}}%
\regfield{RMT\_CH0\_TX\_THR\_EVENT\_INT\_RAW}{1}{12}{{0}}%
\regfield{RMT\_CH3\_ERR\_INT\_RAW}{1}{11}{{0}}%
\regfield{RMT\_CH3\_RX\_END\_INT\_RAW}{1}{10}{{0}}%
\regfield{RMT\_CH3\_TX\_END\_INT\_RAW}{1}{9}{{0}}%
\regfield{RMT\_CH2\_ERR\_INT\_RAW}{1}{8}{{0}}%
\regfield{RMT\_CH2\_RX\_END\_INT\_RAW}{1}{7}{{0}}%
\regfield{RMT\_CH2\_TX\_END\_INT\_RAW}{1}{6}{{0}}%
\regfield{RMT\_CH1\_ERR\_INT\_RAW}{1}{5}{{0}}%
\regfield{RMT\_CH1\_RX\_END\_INT\_RAW}{1}{4}{{0}}%
\regfield{RMT\_CH1\_TX\_END\_INT\_RAW}{1}{3}{{0}}%
\regfield{RMT\_CH0\_ERR\_INT\_RAW}{1}{2}{{0}}%
\regfield{RMT\_CH0\_RX\_END\_INT\_RAW}{1}{1}{{0}}%
\regfield{RMT\_CH0\_TX\_END\_INT\_RAW}{1}{0}{{0}}%
\reglabel{Reset}\regnewline%
\begin{regdesc}\begin{reglist}
\label{fielddesc:RMTREGCHNTXENDINTRAW}\item [RMT\_CH\regindex{n}\_TX\_END\_INT\_RAW] 通道 \regindex{n} 的原始中断状态位。发送结束即触发该中断。（只读）
\label{fielddesc:RMTREGCHNRXENDINTRAW}\item [RMT\_CH\regindex{n}\_RX\_END\_INT\_RAW] 通道 \regindex{n} 的原始中断状态位。接收结束即触发该中断。（只读）
\label{fielddesc:RMTREGCHNERRINTRAW}\item [RMT\_CH\regindex{n}\_ERR\_INT\_RAW] 通道 \regindex{n} 的原始中断状态位。发生错误即触发该中断。（只读）
\label{fielddesc:RMTREGCHNTXTHREVENTINTRAW}\item [RMT\_CH\regindex{n}\_TX\_THR\_EVENT\_INT\_RAW] 通道 \regindex{n} 的原始中断状态位。发送的数据量大于预先设定的值时即触发该中断。（只读）
\label{fielddesc:RMTREGCHNTXLOOPINTRAW}\item [RMT\_CH\regindex{n}\_TX\_LOOP\_INT\_RAW] 通道 \regindex{n} 的原始中断状态位。循环次数达到预先设定的阈值时即触发该中断。（只读）
\end{reglist}\end{regdesc}
\end{register}


\begin{register}{H}{RMT\_INT\_ST\_REG}{\hexprefix{}0054}\label{regdesc:RMTINTSTREG}
\regfield{(reserved)}{12}{20}{000000000000}%
\regfield{RMT\_CH3\_TX\_LOOP\_INT\_ST}{1}{19}{{0}}%
\regfield{RMT\_CH2\_TX\_LOOP\_INT\_ST}{1}{18}{{0}}%
\regfield{RMT\_CH1\_TX\_LOOP\_INT\_ST}{1}{17}{{0}}%
\regfield{RMT\_CH0\_TX\_LOOP\_INT\_ST}{1}{16}{{0}}%
\regfield{RMT\_CH3\_TX\_THR\_EVENT\_INT\_ST}{1}{15}{{0}}%
\regfield{RMT\_CH2\_TX\_THR\_EVENT\_INT\_ST}{1}{14}{{0}}%
\regfield{RMT\_CH1\_TX\_THR\_EVENT\_INT\_ST}{1}{13}{{0}}%
\regfield{RMT\_CH0\_TX\_THR\_EVENT\_INT\_ST}{1}{12}{{0}}%
\regfield{RMT\_CH3\_ERR\_INT\_ST}{1}{11}{{0}}%
\regfield{RMT\_CH3\_RX\_END\_INT\_ST}{1}{10}{{0}}%
\regfield{RMT\_CH3\_TX\_END\_INT\_ST}{1}{9}{{0}}%
\regfield{RMT\_CH2\_ERR\_INT\_ST}{1}{8}{{0}}%
\regfield{RMT\_CH2\_RX\_END\_INT\_ST}{1}{7}{{0}}%
\regfield{RMT\_CH2\_TX\_END\_INT\_ST}{1}{6}{{0}}%
\regfield{RMT\_CH1\_ERR\_INT\_ST}{1}{5}{{0}}%
\regfield{RMT\_CH1\_RX\_END\_INT\_ST}{1}{4}{{0}}%
\regfield{RMT\_CH1\_TX\_END\_INT\_ST}{1}{3}{{0}}%
\regfield{RMT\_CH0\_ERR\_INT\_ST}{1}{2}{{0}}%
\regfield{RMT\_CH0\_RX\_END\_INT\_ST}{1}{1}{{0}}%
\regfield{RMT\_CH0\_TX\_END\_INT\_ST}{1}{0}{{0}}%
\reglabel{Reset}\regnewline%
\begin{regdesc}\begin{reglist}
\label{fielddesc:RMTREGCHNTXENDINTST}\item [RMT\_CH\regindex{n}\_TX\_END\_INT\_ST] \hyperref[int:RMTCHnTXENDINT]{RMT\_CH\regindex{n}\_TX\_END\_INT} 的隐蔽中断状态位。（只读）
\label{fielddesc:RMTREGCHNRXENDINTST}\item [RMT\_CH\regindex{n}\_RX\_END\_INT\_ST] \hyperref[int:RMTCHnRXENDINT]{RMT\_CH\regindex{n}\_RX\_END\_INT} 的隐蔽中断状态位。（只读）
\label{fielddesc:RMTREGCHNERRINTST}\item [RMT\_CH\regindex{n}\_ERR\_INT\_ST] \hyperref[int:RMTCHnERRINT]{RMT\_CH\regindex{n}\_ERR\_INT} 的隐蔽中断状态位。（只读）
\label{fielddesc:RMTREGCHNTXTHREVENTINTST}\item [RMT\_CH\regindex{n}\_TX\_THR\_EVENT\_INT\_ST] \hyperref[int:RMTCHnTXTHREVENTINT]{RMT\_CH\regindex{n}\_TX\_THR\_EVENT\_INT} 的隐蔽中断状态位。（只读）
\label{fielddesc:RMTREGCHNTXLOOPINTST}\item [RMT\_CH\regindex{n}\_TX\_LOOP\_INT\_ST] \hyperref[int:RMTCHnTXLOOPINT]{RMT\_CH\regindex{n}\_TX\_LOOP\_INT} 的隐蔽中断状态位。（只读）
\end{reglist}\end{regdesc}
\end{register}


\begin{register}{H}{RMT\_INT\_ENA\_REG}{\hexprefix{}0058}\label{regdesc:RMTINTENAREG}
\regfield{(reserved)}{12}{20}{000000000000}%
\regfield{RMT\_CH3\_TX\_LOOP\_INT\_ENA}{1}{19}{{0}}%
\regfield{RMT\_CH2\_TX\_LOOP\_INT\_ENA}{1}{18}{{0}}%
\regfield{RMT\_CH1\_TX\_LOOP\_INT\_ENA}{1}{17}{{0}}%
\regfield{RMT\_CH0\_TX\_LOOP\_INT\_ENA}{1}{16}{{0}}%
\regfield{RMT\_CH3\_TX\_THR\_EVENT\_INT\_ENA}{1}{15}{{0}}%
\regfield{RMT\_CH2\_TX\_THR\_EVENT\_INT\_ENA}{1}{14}{{0}}%
\regfield{RMT\_CH1\_TX\_THR\_EVENT\_INT\_ENA}{1}{13}{{0}}%
\regfield{RMT\_CH0\_TX\_THR\_EVENT\_INT\_ENA}{1}{12}{{0}}%
\regfield{RMT\_CH3\_ERR\_INT\_ENA}{1}{11}{{0}}%
\regfield{RMT\_CH3\_RX\_END\_INT\_ENA}{1}{10}{{0}}%
\regfield{RMT\_CH3\_TX\_END\_INT\_ENA}{1}{9}{{0}}%
\regfield{RMT\_CH2\_ERR\_INT\_ENA}{1}{8}{{0}}%
\regfield{RMT\_CH2\_RX\_END\_INT\_ENA}{1}{7}{{0}}%
\regfield{RMT\_CH2\_TX\_END\_INT\_ENA}{1}{6}{{0}}%
\regfield{RMT\_CH1\_ERR\_INT\_ENA}{1}{5}{{0}}%
\regfield{RMT\_CH1\_RX\_END\_INT\_ENA}{1}{4}{{0}}%
\regfield{RMT\_CH1\_TX\_END\_INT\_ENA}{1}{3}{{0}}%
\regfield{RMT\_CH0\_ERR\_INT\_ENA}{1}{2}{{0}}%
\regfield{RMT\_CH0\_RX\_END\_INT\_ENA}{1}{1}{{0}}%
\regfield{RMT\_CH0\_TX\_END\_INT\_ENA}{1}{0}{{0}}%
\reglabel{Reset}\regnewline%
\begin{regdesc}\begin{reglist}
\label{fielddesc:RMTREGCHNTXENDINTENA}\item [RMT\_CH\regindex{n}\_TX\_END\_INT\_ENA] \hyperref[int:RMTCHnTXENDINT]{RMT\_CH\regindex{n}\_TX\_END\_INT} 中断使能位。（读/写）
\label{fielddesc:RMTREGCHNRXENDINTENA}\item [RMT\_CH\regindex{n}\_RX\_END\_INT\_ENA] \hyperref[int:RMTCHnRXENDINT]{RMT\_CH\regindex{n}\_RX\_END\_INT} 中断使能位。（读/写）
\label{fielddesc:RMTREGCHNERRINTENA}\item [RMT\_CH\regindex{n}\_ERR\_INT\_ENA] \hyperref[int:RMTCHnERRINT]{RMT\_CH\regindex{n}\_ERR\_INT} 中断使能位。（读/写）
\label{fielddesc:RMTREGCHNTXTHREVENTINTENA}\item [RMT\_CH\regindex{n}\_TX\_THR\_EVENT\_INT\_ENA] \hyperref[int:RMTCHnTXTHREVENTINT]{RMT\_CH\regindex{n}\_TX\_THR\_EVENT\_INT} 中断使能位。（读/写）
\label{fielddesc:RMTREGCHNTXLOOPINTENA}\item [RMT\_CH\regindex{n}\_TX\_LOOP\_INT\_ENA] \hyperref[int:RMTCHnTXLOOPINT]{RMT\_CH\regindex{n}\_TX\_LOOP\_INT} 中断使能位。（读/写）
\end{reglist}\end{regdesc}
\end{register}


\begin{register}{H}{RMT\_INT\_CLR\_REG}{\hexprefix{}005C}\label{regdesc:RMTINTCLRREG}
\regfield{(reserved)}{12}{20}{000000000000}%
\regfield{RMT\_CH3\_TX\_LOOP\_INT\_CLR}{1}{19}{{0}}%
\regfield{RMT\_CH2\_TX\_LOOP\_INT\_CLR}{1}{18}{{0}}%
\regfield{RMT\_CH1\_TX\_LOOP\_INT\_CLR}{1}{17}{{0}}%
\regfield{RMT\_CH0\_TX\_LOOP\_INT\_CLR}{1}{16}{{0}}%
\regfield{RMT\_CH3\_TX\_THR\_EVENT\_INT\_CLR}{1}{15}{{0}}%
\regfield{RMT\_CH2\_TX\_THR\_EVENT\_INT\_CLR}{1}{14}{{0}}%
\regfield{RMT\_CH1\_TX\_THR\_EVENT\_INT\_CLR}{1}{13}{{0}}%
\regfield{RMT\_CH0\_TX\_THR\_EVENT\_INT\_CLR}{1}{12}{{0}}%
\regfield{RMT\_CH3\_ERR\_INT\_CLR}{1}{11}{{0}}%
\regfield{RMT\_CH3\_RX\_END\_INT\_CLR}{1}{10}{{0}}%
\regfield{RMT\_CH3\_TX\_END\_INT\_CLR}{1}{9}{{0}}%
\regfield{RMT\_CH2\_ERR\_INT\_CLR}{1}{8}{{0}}%
\regfield{RMT\_CH2\_RX\_END\_INT\_CLR}{1}{7}{{0}}%
\regfield{RMT\_CH2\_TX\_END\_INT\_CLR}{1}{6}{{0}}%
\regfield{RMT\_CH1\_ERR\_INT\_CLR}{1}{5}{{0}}%
\regfield{RMT\_CH1\_RX\_END\_INT\_CLR}{1}{4}{{0}}%
\regfield{RMT\_CH1\_TX\_END\_INT\_CLR}{1}{3}{{0}}%
\regfield{RMT\_CH0\_ERR\_INT\_CLR}{1}{2}{{0}}%
\regfield{RMT\_CH0\_RX\_END\_INT\_CLR}{1}{1}{{0}}%
\regfield{RMT\_CH0\_TX\_END\_INT\_CLR}{1}{0}{{0}}%
\reglabel{Reset}\regnewline%
\begin{regdesc}\begin{reglist}
\label{fielddesc:RMTREGCHNTXENDINTCLR}\item [RMT\_CH\regindex{n}\_TX\_END\_INT\_CLR] 用于清除 \hyperref[int:RMTCHnTXENDINT]{RMT\_CH\regindex{n}\_TX\_END\_INT} 中断。（只写）
\label{fielddesc:RMTREGCHNRXENDINTCLR}\item [RMT\_CH\regindex{n}\_RX\_END\_INT\_CLR] 用于清除 \hyperref[int:RMTCHnRXENDINT]{RMT\_CH\regindex{n}\_RX\_END\_INT} 中断。（只写）
\label{fielddesc:RMTREGCHNERRINTCLR}\item [RMT\_CH\regindex{n}\_ERR\_INT\_CLR] 用于清除 \hyperref[int:RMTCHnERRINT]{RMT\_CH\regindex{n}\_ERR\_INT} 中断。（只写）
\label{fielddesc:RMTREGCHNTXTHREVENTINTCLR}\item [RMT\_CH\regindex{n}\_TX\_THR\_EVENT\_INT\_CLR] 用于清除 \hyperref[int:RMTCHnTXTHREVENTINT]{RMT\_CH\regindex{n}\_TX\_THR\_EVENT\_INT} 中断。（只写）
\label{fielddesc:RMTREGCHNTXLOOPINTCLR}\item [RMT\_CH\regindex{n}\_TX\_LOOP\_INT\_CLR] 用于清除 \hyperref[int:RMTCHnTXLOOPINT]{RMT\_CH\regindex{n}\_TX\_LOOP\_INT} 中断。（只写）
\end{reglist}\end{regdesc}
\end{register}




\end{document}
